% =========================================================
% SECCIÓN 3.3 - VARIABLES DE ESTADO Y DOMINIO DEL SISTEMA
% =========================================================

La formulación del modelo DEB a escala de reactor batch exige la definición explícita de las variables de estado, sus unidades físicas, los valores iniciales derivados del protocolo experimental y el dominio matemático que garantiza la existencia de soluciones biológicamente significativas. Esta sección establece el espacio de fases sobre el cual opera el sistema de ecuaciones diferenciales presentado en la Sección~\ref{sec:cap3_formulacion_ode}.

% ---------------------------------------------------------
\subsection{Vector de estado y notación}
\label{subsec:cap3_vector_estado}
% ---------------------------------------------------------

Se define el vector de estado extensivo del reactor $\mathbf{x}(t) \in \mathbb{R}_{\geq 0}^{5}$ como:
\begin{equation}
    \mathbf{x}(t) = \begin{bmatrix} S(t) & E(t) & B(t) & L(t) & C_{\ch{CO2}}(t) \end{bmatrix}^{\!\top},
    \label{eq:cap3_vector_estado}
\end{equation}
donde todas las componentes se expresan en gramos de carbono total del reactor (\si{\gram C}), salvo $S(t)$ que es una cantidad compartida y $E, B, L, C_{\ch{CO2}}$ que resultan del escalamiento de las variables individuales por el número de larvas. La Tabla~\ref{tab:cap3_variables_estado} resume la descripción física de cada variable.

\begin{table}[htbp]
    \centering
    \caption{Variables de estado del modelo DEB escalado a reactor batch.}
    \label{tab:cap3_variables_estado}
    \begin{tabular}{@{}clcc@{}}
        \toprule
        \textbf{Símbolo} & \textbf{Descripción} & \textbf{Unidad} & \textbf{Rango típico} \\
        \midrule
        $S(t)$ & Sustrato orgánico disponible (total reactor) & \si{\gram C} & $[0, S_0]$ \\
        $E(t)$ & Reserva de asimilados (total reactor) & \si{\gram C} & $[0, E_{\max}]$ \\
        $B(t)$ & Biomasa estructural (total reactor) & \si{\gram C} & $[B_{\min}, B_{\max}]$ \\
        $L(t)$ & Lípidos de reserva (total reactor) & \si{\gram C} & $[0, L_{\max}]$ \\
        $C_{\ch{CO2}}(t)$ & Carbono emitido como \ch{CO2} (total reactor) & \si{\gram C} & $[0, C_{\mathrm{tot}}]$ \\
        \bottomrule
    \end{tabular}
\end{table}

La variable $S(t)$ representa la masa de carbono orgánico en el reactor que es físicamente accesible y químicamente degradable por las larvas. Excluye la fracción lignocelulósica recalcitrante, que se agrupa en un término de ``sustrato inerte'' $S_{\mathrm{inerte}}$ constante durante el ciclo. La suma $S(t) + S_{\mathrm{inerte}}$ más el carbono transferido a biomasa, lípidos y \ch{CO2} debe igualar el carbono total inicial $C_{\mathrm{tot}}$, propiedad que se verifica en la Sección~\ref{sec:cap3_analisis}.

La reserva $E(t)$ actúa como amortiguador metabólico: cuando la ingesta excede la demanda combinada de mantenimiento y síntesis, $E$ se acumula; cuando la ingesta es insuficiente, $E$ se moviliza para cubrir el déficit. Esta dinámica de \emph{buffer} evita que la biomasa estructural $B$ se catabolice ante fluctuaciones transitorias de alimento, un mecanismo documentado en el DEB estándar y conservado en la parametrización de Eriksen para BSF \parencite{eriksenDynamicModellingFeed2022}.

La biomasa estructural $B(t)$ se vincula con la biomasa húmeda medida destructivamente mediante el factor de conversión $\rho_{\mathrm{conv}}$ (\si{\gram C \per \gram\;biomasa\;húmeda}). Este parámetro, que depende del contenido de agua y cenizas del tejido larval, se deja como variable libre de calibración en el Capítulo~\ref{chap:hibrido}. La semilla teórica del factor de conversión se sitúa entre $0{,}10$ y $0{,}11$~\si{\gram C \per \gram\;húmeda}, basada en la composición proximal reportada por \textcite{nayakHermetiaIllucensDiptera2023}.

Los lípidos de reserva $L(t)$ modelan el almacenamiento energético previo a la metamorfosis. A diferencia de $B$, que se asume irreversible durante el ciclo larval (sin removilización estructural), $L$ puede experimentar tasas netas negativas si la larva moviliza grasas para mantenimiento en condiciones de ayuno. Sin embargo, en el régimen de alimentación \emph{ad libitum} del protocolo experimental (Sección~\ref{sec:cap2_biomassa}), se espera que $L(t) \geq 0$ durante la mayor parte del ciclo.

La variable $C_{\ch{CO2}}(t)$ acumula el carbono total emitido como dióxido de carbono. Su derivada temporal $\dot{C}_{\ch{CO2}}(t)$ se relaciona directamente con el flujo másico medido por el sensor NDIR mediante la ecuación~\eqref{eq:cap3_puente_sensor}, permitiendo la confrontación modelo--sensor en el Capítulo~\ref{chap:hibrido}.

% ---------------------------------------------------------
\subsection{Variables ambientales y forzamiento externo}
\label{subsec:cap3_forzamiento}
% ---------------------------------------------------------

Además del vector de estado $\mathbf{x}(t)$, el sistema recibe dos variables ambientales medidas por los sensores SHT30 desplegados en el reactor (Sección~\ref{sec:cap2_sensores_ambientales}):
\begin{align}
    T(t) &: \text{temperatura del aire de cabeza libre o del sustrato, en } ^\circ\mathrm{C}, \\
    \mathrm{HR}(t) &: \text{humedad relativa del aire de cabeza libre, en \%}.
\end{align}

Estas variables actúan como forzamiento externo que modula la cinética de asimilación y la tasa metabólica mediante la función de dependencia térmica $f(T)$ y el efecto de la humedad sobre la disponibilidad de sustrato, detallados en la Sección~\ref{sec:cap3_cinetica}. No son variables de estado en el sentido dinámico del sistema, pues su evolución temporal está dictada por el control ambiental de la cámara termohigrométrica ($30^{\circ}\mathrm{C} \pm 2^{\circ}\mathrm{C}$, HR $60$--$80\%$) y no por retroalimentación del metabolismo larval.

Se incluye además una variable pseudo-observable $\hat{B}(t)$, estimada externamente por el calibrador neuronal del Capítulo~\ref{chap:hibrido}, que sirve como entrada dinámica al modelo DEB cuando las mediciones destructivas de biomasa no están disponibles. En la fase de calibración, $\hat{B}(t)$ se reemplaza por el valor interpolado de las mediciones destructivas cada 48~h.

% ---------------------------------------------------------
\subsection{Dominio físico y restricciones de no negatividad}
\label{subsec:cap3_dominio}
% ---------------------------------------------------------

El dominio matemático $\Omega \subset \mathbb{R}_{\geq 0}^{5}$ se define como:
\begin{equation}
    \Omega = \Bigl\{ \mathbf{x} \in \mathbb{R}^{5} \;\Big|\; S \geq 0,\; E \geq 0,\; B \geq B_{\min},\; L \geq 0,\; C_{\ch{CO2}} \geq 0 \Bigr\},
    \label{eq:cap3_dominio}
\end{equation}
donde $B_{\min} > 0$ representa la biomasa estructural mínima compatible con la viabilidad celular. En la práctica experimental la mortalidad larval es despreciable (inferior al $5\%$) en condiciones óptimas \parencite{salamEffectDifferentEnvironmental2022}, por lo que $B_{\min}$ se identifica con la biomasa estructural inicial $B_0$.

Las condiciones iniciales $\mathbf{x}(0) = \mathbf{x}_0$ se derivan directamente del protocolo del Capítulo~\ref{chap:instrumentacion}:
\begin{align}
    S(0) &= S_0 = M_{\mathrm{sustrato}} \cdot w_{\ch{C}} \cdot (1 - f_{\mathrm{inerte}}), \label{eq:cap3_S0} \\
    E(0) &= E_0 \approx 0{,}05 \cdot B_0, \label{eq:cap3_E0} \\
    B(0) &= B_0 = m_{\mathrm{larvas}} \cdot N_{\mathrm{larvas}} \cdot \rho_{\mathrm{conv}}, \label{eq:cap3_B0} \\
    L(0) &= L_0 \approx 0, \label{eq:cap3_L0} \\
    C_{\ch{CO2}}(0) &= C_0 = 0, \label{eq:cap3_C0}
\end{align}
donde $M_{\mathrm{sustrato}} = \SI{250}{\gram}$ es la masa húmeda inicial de sustrato por reactor, $w_{\ch{C}}$ la fracción de carbono en base húmeda del sustrato ($0{,}35$--$0{,}45$ para pulpa de café y cáscaras de frutas), $f_{\mathrm{inerte}}$ la fracción lignocelulósica ($0{,}10$--$0{,}15$), $m_{\mathrm{larvas}} = \SI{0,005}{\gram}$ la biomasa húmeda inicial por larva en estadio L3 sincronizado, y $N_{\mathrm{larvas}} \approx 700$ el número total de larvas sembradas por reactor (Sección~\ref{sec:cap2_biomassa}).

Con estos valores, el balance de carbono inicial se aproxima como:
\begin{align*}
    S_0 &\approx \SI{250}{\gram} \times 0{,}40 \times 0{,}85 \approx \SI{85}{\gram C}, \\
    B_0 &\approx 700 \times \SI{0,005}{\gram} \times 0{,}10 \approx \SI{0,35}{\gram C}, \\
    E_0 &\approx 0{,}05 \times \SI{0,35}{\gram C} \approx \SI{0,018}{\gram C}, \\
    L_0 &\approx 0, \qquad C_0 = 0,
\end{align*}
de modo que $C_{\mathrm{tot}} = S_0 + S_{\mathrm{inerte}} + B_0 + E_0 \approx \SI{100}{\gram C}$. La reserva inicial $E_0$ se fija en un $5\%$ de la biomasa estructural inicial, valor consistente con larvas recién transferidas del medio de cría al reactor experimental, donde la ingesta previa ha sido limitada. La condición $L_0 = 0$ refleja que los estadios tempranos (L3 inicial) poseen reservas lipídicas mínimas; la acumulación de $L$ se activa en estadios posteriores mediante la función de prioridad de asignación descrita en la Sección~\ref{subsec:cap3_flujos}.

% ---------------------------------------------------------
\subsection{Escalamiento de individuo a reactor batch}
\label{subsec:cap3_escalamiento}
% ---------------------------------------------------------

El modelo DEB de \textcite{eriksenDynamicModellingFeed2022} está formulado para una larva individual, con variables expresadas por organismo (\si{\gram C \per larva}). Para adaptarlo al reactor batch, se adopta la hipótesis de equivalencia fisiológica: se asume que todas las larvas del reactor comparten el mismo estado metabólico promedio, de modo que las variables de estado a escala de reactor son el producto del número de larvas $N_{\mathrm{larvas}}$ por la variable individual.

Formalmente, si $\tilde{e}$, $\tilde{b}$, $\tilde{l}$ y $\tilde{c}_{\ch{CO2}}$ denotan las variables por larva, entonces:
\begin{equation}
    E = N_{\mathrm{larvas}} \cdot \tilde{e}, \quad
    B = N_{\mathrm{larvas}} \cdot \tilde{b}, \quad
    L = N_{\mathrm{larvas}} \cdot \tilde{l}, \quad
    C_{\ch{CO2}} = N_{\mathrm{larvas}} \cdot \tilde{c}_{\ch{CO2}}.
    \label{eq:cap3_escalamiento}
\end{equation}

Esta hipótesis es razonable cuando la densidad larval se mantiene dentro del rango operativo del protocolo experimental ($\sim$\SI{4}{\gram\per\litre} de biomasa húmeda en \SI{0,8}{\litre}), donde la competencia por el alimento es limitada y la heterogeneidad de crecimiento entre individuos es baja \parencite{dienerConversionOrganicMaterial2011,guillaumeAsymptoticEstimatedDigestibility2023}. En densidades superiores, la variabilidad interindividual aumenta y el escalamiento medio puede subestimar la dispersión de la respiración; esta limitación se mitiga en el Capítulo~\ref{chap:hibrido} mediante la incertidumbre del estimador neuronal.

El número de larvas $N_{\mathrm{larvas}}$ se trata como constante durante el ciclo de 12--15 días, dado que la mortalidad registrada en condiciones controladas es inferior al $5\%$ \parencite{salamEffectDifferentEnvironmental2022}. Cualquier reducción por mortalidad se absorbe en la calibración del parámetro de eficiencia de conversión $\kappa_G$.

En resumen, esta sección ha definido el espacio de estados, las condiciones iniciales experimentales y el dominio matemático del modelo. El siguiente paso consiste en formular las ecuaciones diferenciales que gobiernan la evolución de $\mathbf{x}(t)$ dentro de $\Omega$, tema de la Sección~\ref{sec:cap3_formulacion_ode}.